\subsection{ABB Mobile YuMi}
\paragraph{TODO 1} Rewrite SMCYuMiNode into RealMobileYuMiRobotManagerNode (same way as HeronNode)
\paragraph{TODO 2} Finish the todos in yuminode, delete commented code

\subsubsection{Hardware}
YuMi has 2 7dof arms. The payload on the end-effector is like 0.5kg (super weak).
The mechanical design is such that the first 3 links of the arms are very close
to the torso, so it's very easy to get self-collisions between the elbow and the torso.
In short, the elbows should be on the side of the robot,
and the end-effectors should be in front in almost all cases
(have the range of a person with limited shoulder mobility).

The in-house developed research-level mobile base
is omnidirectional and generally works quite well, 
but with some quirks. Namely,
when the robot stops, the wheels reorient themselves
such that the robot can not be pushed.
A similar kind of in-place re-orientation can happen 
if there is a large change in the angle of the commanded velocity.
These functions are performed by the platform - one commands
the base velocity in twist-form. Direct wheel commands are inaccessible.
Additionally, the base only acts of velocities of a certain minimum norm.
Velocities lower than that are (mostly) discarded.
In total, this means that the motion of base should be a trajectory
of (relatively) low curvature, and traversed at a speed
over the minimum. In other words, in continuous operation the robot
is closer to a differential drive. The omnidirectionality is useful
for going in an arbitrary direction from a stationary position.
TODO: None of this is explicitly modelled in YuMiRobotManager,
and it really should be for experiments to be successfully designed
(you can't eyeball your way around this, it has to be modelled).


\subsubsection{Software on the robot}
The robot has a full ROS2 stack.
ROS connects to something called EGM which
is a common external-control controller, i.e. it just listens
to incoming velocity commands and then executes
them with a proprietary inner controller. This is referred to as EGM mode on the robot.
EGM is super annoying (randomly stops working, so need to reset it 
etc). The ROS driver is equally as annoying. Avoid this robot if you have other options,
unless you like restarting everything (5ish min process) every half hour.

The robot has an $f_{ \text{ctrl} } \approx  $60Hz interface for full-body velocity control,
which is a combination of individual velocity ``controllers'' for the base, 
the torso pillar (which we don't use) and the mobile base.

For the most part, the communication with the robot goes over 5G WiFi in the lab,
meaning no code needs to be installed on the robot, but there are occasional lag spikes
leading to loss of control for (eyeballed) 0.2-0.5 seconds. 
[Marko's estimate is that ] This kind of interfacing is just barely acceptable
if the control is slow and conservative. It is possible to install additional ROS2
packages on the robot, but [by Marko's estimate] installing the control software stack
is an ill-advised adventure as there is no Docker on the robot.

From a control perspective,
the specific topics to interact with the robot are:
\begin{itemize}
		\item (namespace)/platform/joint\_states - used
				to fetch current joint angles (all of them).
				Uses sensor\_msgs.msg.JointState as message format,
				and it's just a vector of all joint angles.
		\item (namespace)/platform/joints\_cmd - used
				to send a velocity command to all 
				joints. Also uses  sensor\_msgs.msg.JointState
		\item (namespace)/platform/odometry - odometry topic.
				Good enough for small experiments,
				but there is AMCL for actual localization
				(0.5Hz though, you need a filter to make use of
				it).
\end{itemize}

\subsubsection{Concrete steps to use the robot}
TODO document how to operate mobile yumi in vasteros 
(sequence of buttons to press to restart it after it craps out),
how to start Rviz and what to click in it to start EGM etc.

\paragraph{In code}
ATM in python/smc/multiprocessing/smc\_yumi\_node.py,
will get rewriten to be like heron
and moved to python/smc/robots/implementations/mobile\_yumi.py

\subsection{Sleipner + YuMi (SLuMi)}
TODO
It's ROS1
things to do
1. turn power on on robot. first 
